% ENPM703 Group 2 report, set in the CVPR 2017 author-kit style.
% The preamble follows the kit's egpaper_final.tex; cvpr.sty and ieee.bst are
% the kit's files, unmodified.
%
% Build:  latexmk -pdf report.tex

\documentclass[10pt,twocolumn,letterpaper]{article}

\usepackage{cvpr}
\usepackage{times}
\usepackage{epsfig}
\usepackage{graphicx}
\usepackage{amsmath}
\usepackage{amssymb}
\usepackage{xcolor}

% Include other packages here, before hyperref.

\graphicspath{{figures/}}

% Visible placeholder. Delete each one as its section is written; a finished
% report has no \todo left in it.
\newcommand{\todo}[1]{\textcolor{red}{[TODO: #1]}}

% If you comment hyperref and then uncomment it, delete report.aux before
% re-running latex.
\usepackage[breaklinks=true,bookmarks=false]{hyperref}

% Final-copy layout: named authors, no review ruler. Comment this line out to
% get the anonymous review layout with line numbers.
\cvprfinalcopy

\def\cvprPaperID{****} % printed in the review layout only

\begin{document}

%%%%%%%%% TITLE
\title{Repository-Aware Code Completion under a Context Budget}

\author{Marcus Hurt \quad Maria Vilanova Garcia \quad Kenneth Sutter \quad Aditi Thangavelu\\
University of Maryland, College Park\\
ENPM703, Group 2
}

\maketitle

%%%%%%%%% ABSTRACT
\begin{abstract}
   \todo{One paragraph: the problem, what we varied, the data and model, and
   the headline result.}
\end{abstract}

%%%%%%%%% BODY TEXT
\section{Introduction}

\todo{Why cross-file context matters for code completion, why the context
budget forces a choice, and the two questions this project asks: what should
fill a budget of $N$ tokens for a small frozen code model, and can a cheap
classifier predict the best choice per completion.}

%-------------------------------------------------------------------------
\section{Related Work}

RepoFuse~\cite{liang2024repofuse} compares signature context against
retrieved-snippet context at matched truncation and reports that fusing the
two beats either alone. Repoformer~\cite{wu2024repoformer} learns per example
whether to retrieve at all, and ProCC~\cite{tan2024procc} uses a contextual
bandit to pick a retrieval strategy per example.

\todo{Expand, and position our two passes against these.}

%-------------------------------------------------------------------------
\section{Method}

\todo{The context builder: the methods compared, how each fills the token
budget, and the prompt layout. Then the pass 2 selector.}

%-------------------------------------------------------------------------
\section{Experiments}

\todo{Setup: RepoBench~\cite{liu2023repobench} Python as the development set
and CrossCodeEval~\cite{ding2023crosscodeeval} Python as the final test,
DeepSeek-Coder-1.3B~\cite{guo2024deepseekcoder} in 4-bit, the budgets swept,
the metrics (exact match, edit similarity, completion time), and the paired
statistical tests.}

\todo{Results and discussion.}

% Figure and table patterns from the author kit. Size a figure as a multiple
% of \linewidth, and keep its text readable in a printed copy.
%
% \begin{figure}[t]
% \begin{center}
%    \includegraphics[width=0.8\linewidth]{myfigure.pdf}
% \end{center}
%    \caption{One-column figure. Use figure* for a figure spanning both columns.}
% \label{fig:example}
% \end{figure}
%
% \begin{table}
% \begin{center}
% \begin{tabular}{|l|c|c|}
% \hline
% Method & EM & ES \\
% \hline\hline
% Method A & -- & -- \\
% Method B & -- & -- \\
% \hline
% \end{tabular}
% \end{center}
% \caption{Results.}
% \label{tab:example}
% \end{table}

%-------------------------------------------------------------------------
\section{Conclusion}

\todo{What the results say about each question, the limits of the study, and
what we would do next.}

%-------------------------------------------------------------------------
\section*{Contributions}

\todo{Who did what.}

\section*{AI and external-code use}

\todo{Summarize from AI\_USE.md in the repository; the course requires
attribution for LLM-generated and externally sourced code and text.}

{\small
\bibliographystyle{ieee}
\bibliography{references}
}

\end{document}
